\begingroup
\renewcommand{\arraystretch}{1.45}
\fancyhead[L]{\small LEMNISCATIC FUNCTIONS II.}
\Needspace{18\baselineskip}
\CurvaArticle{12}{[12.]}
\begin{center}Various hidden series summations.\end{center}
\Needspace{14\baselineskip}
\[
\begin{aligned}1^{0}\quad&
\biggl[\dfrac{2}{e^\pi+e^{-\pi}}\biggr]^2
+\biggl[\dfrac{2}{e^{2\pi}+e^{-2\pi}}\biggr]^2
+\biggl[\dfrac{2}{e^{3\pi}+e^{-3\pi}}\biggr]^2+\mathord{.}\mathord{.}\mathord{.}
\\[8pt]&=\dfrac12\dfrac{\varpi\varpi}{\pi\pi}-\dfrac{1}{2\pi}-\dfrac12\end{aligned}
\]
\Needspace{14\baselineskip}
\[
\begin{aligned}2^{0}\quad&
\biggl[\dfrac{2}{e^{\frac12\pi}-e^{-\frac12\pi}}\biggr]^2
+\biggl[\dfrac{2}{e^{\frac32\pi}-e^{-\frac32\pi}}\biggr]^2
+\biggl[\dfrac{2}{e^{\frac52\pi}-e^{-\frac52\pi}}\biggr]^2+\mathord{.}\mathord{.}\mathord{.}
\\[8pt]&=\dfrac12\dfrac{\varpi\varpi}{\pi\pi}-\dfrac{1}{2\pi}\end{aligned}
\]
\Needspace{14\baselineskip}
\[
\begin{aligned}3^{0}\quad&
\biggl[\dfrac{2}{e^\pi-e^{-\pi}}\biggr]^2
+\biggl[\dfrac{2}{e^{2\pi}-e^{-2\pi}}\biggr]^2
+\biggl[\dfrac{2}{e^{3\pi}-e^{-3\pi}}\biggr]^2+\mathord{.}\mathord{.}\mathord{.}
\\[8pt]&=\dfrac16-\dfrac{1}{2\pi}\end{aligned}
\]
\Needspace{14\baselineskip}
\[
\begin{aligned}4^{0}\quad&
\biggl[\dfrac{2}{e^{\frac12\pi}+e^{-\frac12\pi}}\biggr]^2
+\biggl[\dfrac{2}{e^{\frac32\pi}+e^{-\frac32\pi}}\biggr]^2
+\biggl[\dfrac{2}{e^{\frac52\pi}+e^{-\frac52\pi}}\biggr]^2+\mathord{.}\mathord{.}\mathord{.}
\\[8pt]&=\dfrac{1}{2\pi}\end{aligned}
\]
\Needspace{12\baselineskip}
\[
\begin{aligned}5^{0}\quad&
\dfrac{2}{e^{\frac12\pi}+e^{-\frac12\pi}}
+\dfrac{2}{e^{\frac32\pi}+e^{-\frac32\pi}}+\mathord{.}\mathord{.}\mathord{.}
\\[8pt]&=\dfrac12\dfrac\varpi\pi\end{aligned}
\]
\nopagebreak[4]
\[
\begin{aligned}6^{0}\quad&\end{aligned}
\]



\Needspace{18\baselineskip}
\SourceLeaf{430}
\CurvaArticle{13}{[13.]}
\begin{center}\textit{Calculations pertaining to the lemniscate.}\end{center}
\Needspace{18\baselineskip}
\[
\begin{aligned}
\{(\sin\operatorname{lemn}\dfrac25\varpi)^2
+(\sin\operatorname{lemn}\dfrac45\varpi)^2\}^2&=14\surd5-30\\[5pt]
\{(\sin\operatorname{lemn}\dfrac25\varpi)^2
-(\sin\operatorname{lemn}\dfrac45\varpi)^2\}^2&=10\surd5-22
\end{aligned}
\]
\Needspace{16\baselineskip}
\[
\begin{gathered}
\dfrac12(\sin\operatorname{lemn}\dfrac25\varpi)^4
+\dfrac12(\sin\operatorname{lemn}\dfrac45\varpi)^4
\\[5pt]=6\surd5-13\\[5pt]
=0.4164078649\ 9873817845\ 5042012387\ 65741
\end{gathered}
\]
\Needspace{16\baselineskip}
\[
\begin{gathered}
\{\dfrac12(\sin\operatorname{lemn}\dfrac25\varpi)^4
-\dfrac12(\sin\operatorname{lemn}\dfrac45\varpi)^4\}^4
\\[5pt]=340-152\surd5\\[5pt]
=0.1176674200\ 3196614580\ 5602352846\ 01228
\end{gathered}
\]
\Needspace{10\baselineskip}
\noindent hence the root
\[
0.3430268503\ 0761971797\ 7310507555\ 85731
\]
\noindent thus
\Needspace{12\baselineskip}
\[
\begin{aligned}(\sin\operatorname{lemn}\dfrac45\varpi)^4
&=\\[5pt]&\quad{}0.0733810146\ 9111846047\ 7731504831\ 80010\end{aligned}
\]
\Needspace{12\baselineskip}
\[
\begin{aligned}(\sin\operatorname{lemn}\dfrac25\varpi)^4
&=\\[5pt]&\quad{}0.7594347153\ 0635789643\ 2352519943\ 51472\end{aligned}
\]
\Needspace{12\baselineskip}
\[
\begin{aligned}(\sin\operatorname{lemn}\dfrac45\varpi)^2
&=\\[5pt]&\quad{}0.2708893033\ 8999814497\ 30710\end{aligned}
\]
\Needspace{12\baselineskip}
\[
\begin{aligned}(\sin\operatorname{lemn}\dfrac25\varpi)^2
&=\\[5pt]&\quad{}0.8714555153\ 9155336074\ 646029\end{aligned}
\]
\Needspace{18\baselineskip}
\[
\begin{aligned}
Q\dfrac1{10}\varpi
&={}^{5}\!\surd(\dfrac34\cos\dfrac1{10}\pi+\dfrac14\sin\dfrac1{10}\pi)\\[5pt]
Q\dfrac2{10}\varpi
&={}^{5}\!\surd(\dfrac34\cos\dfrac1{10}\pi-\dfrac14\sin\dfrac1{10}\pi)
\end{aligned}
\]
\Needspace{18\baselineskip}
\CurvaArticle{14}{[14.]}
\begin{center}\textit{Lemniscatic function.}\end{center}
\Needspace{10\baselineskip}
\[
\int\dfrac{\mathrm{d}x}{\surd(1-x^4)}=\varphi,
\qquad
\int\dfrac{\mathrm{d}X}{\surd(1-X^4)}=\Phi
\]
\Needspace{18\baselineskip}
\[
\begin{array}{rcccc}
\text{Let}&1&x&\surd(1-xx)&\surd(1+xx)\\[5pt]
\text{be as}&p&q&r&s\\[7pt]
\text{and}&1&X&\surd(1-XX)&\surd(1+XX)\\[5pt]
\text{be as}&P&Q&R&S\quad.
\end{array}
\]
\Needspace{18\baselineskip}
\noindent then the quantities corresponding to $\varphi+\Phi$ are


\SourceLeaf{431}
\begingroup\small\setlength{\arraycolsep}{4pt}\[
\begin{array}{@{}ll|l|l|l@{}}
\text{as}&ppPP+qqQQ&pqRS+rsPQ&prPR-qsQS&psPS+qrQR\\
\text{or}&pqRS-rsPQ&qqPP-ppQQ&qrPS-psQR&qsPR-prQS\\
\text{or}&prPR+qsQS&qrPS+psQR&ppRR-qqSS&2pqPQ+rsRS\\
&&&=rrPP-ssQQ&\\
\text{or}&psPS-qrQR&qsRR+prQS&rsRS-2pqPQ&ssPP+rrQQ\\
&&&&=ppSS+qqRR
\end{array}
\]\endgroup
\Needspace{12\baselineskip}
\noindent Let
\[
\begin{aligned}
\sin\operatorname{lemn}X&=x\\[5pt]
\sin\operatorname{lemn}Y&=y\\[5pt]
\sin\operatorname{lemn}Z&=z
\end{aligned}
\]
\Needspace{12\baselineskip}
\noindent and let $X,\ Y,\ Z$ depend on one another in such a way that
\[
X+Y+Z=0
\]
\Needspace{10\baselineskip}
\noindent Set
\[
\int xx\,\mathrm{d}X=FX
\]
\noindent and
\[
FX+FY+FZ=u
\]
\Needspace{18\baselineskip}
\noindent From formulas encountered earlier it follows that
\[
\begin{aligned}
xx-zz&=yz\surd(1-x^4)-xy\surd(1-z^4)\\[5pt]
yy-zz&=xz\surd(1-y^4)-xy\surd(1-z^4)\\[5pt]
xx-yy&=yz\surd(1-x^4)-xz\surd(1-y^4)
\end{aligned}
\]
\Needspace{16\baselineskip}
\noindent From
\[
\begin{aligned}
\mathrm{d}u&=\dfrac{xx\,\mathrm{d}x}{\surd(1-x^4)}
+\dfrac{yy\,\mathrm{d}y}{\surd(1-y^4)}
+\dfrac{zz\,\mathrm{d}z}{\surd(1-z^4)}\\[5pt]
0&=\dfrac{\mathrm{d}x}{\surd(1-x^4)}
+\dfrac{\mathrm{d}y}{\surd(1-y^4)}
+\dfrac{\mathrm{d}z}{\surd(1-z^4)}
\end{aligned}
\]
\Needspace{22\baselineskip}
\noindent it follows that
\[
\begin{aligned}
\mathrm{d}u&=(yy-xx)\dfrac{\mathrm{d}y}{\surd(1-y^4)}
+(zz-xx)\dfrac{\mathrm{d}z}{\surd(1-z^4)}\\[5pt]
&=xz\,\mathrm{d}y+xy\,\mathrm{d}z-yz.\surd(1-x^4).
\\[7pt]&\quad{}\biggl[\dfrac{\mathrm{d}y}{\surd(1-y^4)}+\dfrac{\mathrm{d}z}{\surd(1-z^4)}\biggr]\\[5pt]
&=xz\,\mathrm{d}y+xy\,\mathrm{d}z+yz\,\mathrm{d}x=\mathrm{d}.xyz
\end{aligned}
\]
\Needspace{12\baselineskip}
\SourceLeaf{432}

thus
\[
u=xyz
\]
Consequently,
\Needspace{12\baselineskip}
\[
\begin{aligned}F(a+b)&=Fa+Fb\\[7pt]&\quad{}-\sin\operatorname{lemn}a.\sin\operatorname{lemn}b.\sin\operatorname{lemn}(a+b)\end{aligned}
\]
If one sets
\[
Fa-\frac{a}{90^0}F90^0=Ga
\]
\Needspace{12\baselineskip}
then
\[
\begin{aligned}G(a+b)&=Ga+Gb\\[7pt]&\quad{}-\sin\operatorname{lemn}a.\sin\operatorname{lemn}b.\sin\operatorname{lemn}(a+b)\end{aligned}
\]
\Needspace{10\baselineskip}
\[
Ga=\frac{\mathrm dQa}{Qa.\mathrm da}
\]

\clearpage
\endgroup
